% ===========================================================================
%  Capítulo 2 — As Leis da Magia
% ===========================================================================
\chapter{As Leis da Magia}

% ---------------------------------------------------------------------------
\section{As Categorias da Magia}

O pilar da Magia tem \textbf{oito escolas}, divididas em dois grupos:

\begin{itemize}
    \item \textbf{Magia Ofensiva (4 escolas):} os quatro elementos clássicos --- Água,
          Fogo, Terra e Vento. É onde mora a maior quantidade de feitiços do livro.
    \item \textbf{Magia de Suporte (4 escolas):} Cura, Desintoxicação, Barreira e
          Invocação. Invocação é a menor lista de feitiços do livro, e a única cujo efeito
          principal age sozinho depois de conjurado.
\end{itemize}

Os dois grupos usam exatamente as mesmas regras: a mesma tabela de PM, o mesmo tempo de
conjuração, o mesmo Bônus de Rank. A divisão é de assunto, não de mecânica.

\begin{nota}{Quão raro é um mago}
    Apenas 1 em cada 20 pessoas nasce com capacidade de manipular mana. Dessas, apenas 1
    em cada 20 consegue treinar o suficiente pra virar mago de verdade --- cerca de 1
    pessoa em 400. De cada cem magos formados, só um completa os estudos até o rank
    Avançado: um mago Avançado é aproximadamente 1 em 40.000 pessoas.
\end{nota}

% ---------------------------------------------------------------------------
\section{A Regra de Encantamentos}

A regra de ouro: o poder da magia depende do encantamento e do tempo gasto para conjurá-la.
O tamanho do encantamento é proporcional ao rank --- quanto maior, mais longo o cântico.

\begin{itemize}
    \item \textbf{Conjuração Padrão:} recita o cântico inteiro. 100\% do efeito.
    \item \textbf{Encantamento Encurtado:} pula versos intencionalmente. Mais rápido, mas
          instável.
    \item \textbf{Conjuração Silenciosa:} manipula a mana diretamente, sem palavra alguma.
          O método mais raro e mais flexível.
\end{itemize}

\begin{aviso}{Penalidade do Encantamento Encurtado}
    Ao encurtar, role \textbf{metade dos dados de dano} (arredondado pra baixo) e a área de
    efeito é reduzida em \textbf{um terço}. O BC continua sendo somado integralmente --- a
    maestria não some, só a estrutura do feitiço fica instável.
\end{aviso}

\begin{nota}{Conjuração Silenciosa --- regra completa, sem exceção escondida}
    Conjurar em silêncio sempre faz exatamente estas quatro coisas ao mesmo tempo, nesta
    ordem de leitura:

    \begin{enumerate}
        \item \textbf{Custo de Ação:} sempre usa a coluna \enquote{Silenciosa} da tabela da
              seção 3 --- sempre mais rápida que a Padrão, e \textbf{nunca mais lenta} que
              a Encurtada. Nos dois primeiros patamares as duas empatam em 1 Ação (e no
              Principiante a Silenciosa ainda ganha a primeira do turno de graça); do
              Avançado em diante a Silenciosa abre vantagem de verdade.
        \item \textbf{Dano:} metade dos dados, arredondado pra baixo --- o mesmo valor do
              Encantamento Encurtado, nunca mais que isso por padrão.
        \item \textbf{Área:} reduzida em um terço --- de novo, o mesmo valor do Encurtado.
        \item \textbf{Bônus de Forma} (sempre incluso, nunca opcional de pagar): escolha
              um, de graça, toda vez que conjurar em silêncio --- dobrar o alcance, mudar o
              formato da área (linha $\leftrightarrow$ cone $\leftrightarrow$ esfera), ou
              segurar o disparo por até 1 turno.
    \end{enumerate}

    \medskip
    \textbf{Quando um talento ou Antecedente diz \enquote{sem sofrer a penalidade de
    dano/área}}, ele remove só o item 2 e/ou o item 3 --- exatamente os que ele nomear. Os
    itens 1 e 4 nunca são removidos por nada, porque eles não são a penalidade: são a
    definição estrutural do próprio método.

    \medskip
    A Conjuração Silenciosa não é comprável com PA --- vem de um Antecedente, de uma raça,
    ou de uma Maestria de Rank alto.
\end{nota}

\begin{aviso}{A única exceção à regra de \enquote{não existe ação bônus}}
    O Capítulo 4, §3 diz que tudo neste sistema é medido em Ações e que \textbf{não existe
    ação bônus}. A Conjuração Silenciosa de rank Principiante é a única exceção nomeada do
    livro: a primeira delas em cada turno é \textbf{gratuita}.

    \medskip
    Ela existe porque, sem ela, conjurar em silêncio num rank baixo custaria uma Ação
    inteira pra entregar metade dos dados e dois terços da área --- ninguém usaria nunca, e
    o método mais característico do mundo morreria na ficha. Vale só pro rank Principiante,
    só pra primeira do turno, e o livro não abre nenhuma outra: se você encontrar qualquer
    outra coisa que se comporte como ação bônus, é erro de texto, não regra.

    \medskip
    \textbf{E ela não conta no Teto de Ações} (Cap. 4, §4: 5 por turno, no máximo 2
    externas). Ela não gasta Ação nenhuma, então não há Ação pra contar --- mas ela também
    não é uma das 2 externas, e um Imperador de magia continua limitado a uma por turno, do
    rank mais fraco que ele conhece.
\end{aviso}

% ---------------------------------------------------------------------------
\section{Tempo de Conjuração por Rank}

A tabela de tempo de conjuração está no Capítulo 1, junto das outras tabelas de Rank ---
ela governa quantas Ações uma magia custa, e o cântico cresce com o rank. Magias grandes
exigem que o grupo proteja o mago enquanto ele canta.

Como o Capítulo 4 permite \termo{dividir o cântico} entre turnos, magias de 4, 5 ou 6 Ações
são perfeitamente jogáveis --- elas só exigem que alguém segure a linha de frente enquanto
o apocalipse é preparado.

% ---------------------------------------------------------------------------
\section{Combinações entre Árvores}

Duas árvores em Rank Avançado ou superior não competem pelo mesmo personagem --- elas se
somam.

\subsection{Magia Combinada}

Magias Combinadas são feitiços compostos por duas ou mais magias conjuradas em sequência,
cujo resultado é maior que a soma das partes. Nova Congelante é literalmente Vento $+$
Água/Gelo; Vapor Seco é Vento $+$ Fogo.

\textbf{Requisito:} rank Avançado ou superior em \textbf{ambas} as escolas envolvidas --- o
mesmo portão que a seção seguinte usa pra qualquer par de árvores. Não existe combinar uma
escola que você domina com outra que você mal começou: a magia composta exige que as duas
metades estejam maduras.

\begin{itemize}
    \item Conjure as duas magias no mesmo turno ou em turnos consecutivos.
    \item Pague o PM e as Ações de ambas.
    \item O resultado é uma terceira magia, com efeito próprio.
\end{itemize}

\begin{tabelalivro}{L{4.6cm} L{3.2cm} L{6cm}}{Magias Combinadas catalogadas}
\cabtabela{\textbf{Combinação} & \textbf{Resultado} & \textbf{Efeito}}
Respingos de Água $+$ Campo de Gelo & Nova Congelante & Todos na área ficam Molhados e imediatamente Congelados, sem teste. Dano 4d8 de frio. \\
Tempestade $+$ qualquer magia de Fogo & Névoa Escaldante & Esfera de 18m de vapor. Escuridão total $+$ 1d6 de dano ígneo por turno a quem estiver dentro. \\
Canhão de Água $+$ Lâmina de Gelo & Serra d'Água & A linha do Canhão passa a causar dano cortante e ignora metade da CA de armaduras não-mágicas. \\
Cumulonimbus $+$ Nevasca & Inverno Rasgado & A tempestade vira granizo. Toda criatura na área de 1,5 km sofre 2d6 de frio por turno, sem teste. \\
Muralha de Pedra $+$ qualquer magia de Fogo & Vidro Cortante & A muralha derrete e resolidifica em lâminas de vidro. Quem tocar ou atravessar sofre 3d8 cortante; a muralha perde metade dos PV Máximos. \\
Lama Viva $+$ Respingos de Água & Areia Movediça Instantânea & Esfera de 9m vira lodo: Deslocamento reduzido à metade e teste de Força (CD 8 $+$ BC) ou fica Atolado até a cintura. \\
Cordilheira $+$ Ciclone & Tempestade de Poeira & Esfera de 18m: Cego pra quem estiver dentro, e projéteis mundanos erram automaticamente até a poeira baixar (1 minuto). \\
\end{tabelalivro}

\begin{nota}{Regra de Ouro para o Mestre}
    Se o jogador descrever uma combinação que faz sentido físico, deixe funcionar e invente
    o efeito na hora. Este sistema recompensa quem pensa como cientista --- foi assim que
    Rudeus criou metade do arsenal dele.
\end{nota}

\subsection{Combinações Além da Magia}

A mesma lógica funciona entre \emph{qualquer} duas árvores em Rank Avançado ou superior,
mesmo cruzando pilares diferentes (Magia $+$ Corpo, Magia $+$ Utilidade, Corpo $+$
Utilidade). O requisito e o custo são os mesmos: Avançado ou superior nas duas árvores
envolvidas, e você paga o custo de cada lado inteiro (Ação, PM, PT ou PP). O efeito nunca é
permanente, a menos que a tabela diga o contrário.

\begin{tabelalivro}{L{4.4cm} L{3cm} L{6.4cm}}{Combinações entre pilares}
\cabtabela{\textbf{Combinação} & \textbf{Resultado} & \textbf{Efeito}}
Deus da Espada $+$ Magia de Fogo & Lâmina em Chamas & Gaste a Ação e o PM de uma magia de Fogo de rank Avançado ou inferior, mais 1 PT: por 1 minuto, seu Dado de Arma causa $+$1d8 de dano ígneo extra. \\
Tático $+$ Magia de Terra & Terreno Escolhido & Antes de um combate previsto, gaste 1 PP pra declarar que já preparou o chão. Se a luta acontecer lá, sua próxima magia de Terra tem $+$50\% de área. \\
Ladino $+$ Magia de Invocação & Familiar Furtivo & Sua próxima invocação nasce com o seu Bônus de Rank de Ladino em Furtividade, e reporta o que viu sem gastar sua Ação. \\
Bardo $+$ Magia de Cura & Canção que Cura & Enquanto sustentar uma Canção, cada magia de Cura que você conjurar recupera $+$1d8 extra em todos os alvos afetados. \\
Cavalaria e Escudos $+$ Barreira & Escudo Vivo & Uma vez por combate, gaste 1 PM: seu escudo ganha PV temporários iguais ao seu Bônus de Rank, absorvidos antes de a CA importar. \\
\end{tabelalivro}

\begin{nota}{Quando a combinação vira uma árvore própria}
    Às vezes duas árvores em Rank Avançado se encaixam bem demais pra caber numa única
    habilidade --- o Estilo Deus do Norte com a Magia de Vento, por exemplo, virou o
    \textbf{Estilo Vendaval}, uma décima oitava sub-árvore inteira que só se revela pra quem
    cumpriu os dois pré-requisitos. Isso não é a regra --- é o teto dela.
\end{nota}

% ---------------------------------------------------------------------------
\section{Maestrias}

Ao desbloquear um Rank em qualquer árvore, você recebe imediatamente e de graça a Maestria
correspondente. Ela não custa PA, não conta como conhecimento, e não pode ser recusada.

A lógica é simples: subir de rank não é decorar mais um feitiço --- é compreender o
elemento de um jeito novo. Roxy não comprou a habilidade de encurtar cânticos; ela entendeu
água fundo o suficiente pra que encurtar virasse natural.

As Maestrias estão listadas dentro de cada árvore, em cada Rank, marcadas com
\simbmaestria{} no Capítulo 3.
